\documentclass[10pt]{article}

%%% Doc layout
\usepackage{fullpage} 
\usepackage{booktabs}       % professional-quality tables
\usepackage{microtype}      % microtypography
\usepackage{parskip}
\usepackage{times}

%% Hyperlinks always black, no weird boxes
\usepackage[hyphens]{url}
\usepackage[colorlinks=true,pdfborder={0 0 0}]{hyperref}

%%% Math typesetting
\usepackage{amsmath,amssymb}

%%% Write out problem statements in purple, solutions in black
\usepackage{xcolor}
\newcommand{\officialdirections}[1]{{\color{purple} #1}}

%%% Avoid automatic section numbers (we'll provide our own)
\setcounter{secnumdepth}{0}

%% --------------
%% Header
%% --------------
\usepackage{fancyhdr}
\fancyhf{}
\setlength{\headheight}{15pt}
\fancyhead[C]{CS 136 - 2026f - PS4 Submission}

\fancyfoot[C]{\thepage} % page number
\renewcommand\headrulewidth{0pt}
\pagestyle{fancy}
\newcommand{\argmax}{\operatorname*{arg\,max}}


%% --------------
%% Begin Document
%% --------------
\begin{document}

~\\ %% add vertical space

{\Large{\bf Student Name: TODO}}

~\\ %% add vertical space
{\bf Collaboration Statement:}

Total hours spent: TODO

I consulted the following human, textbook, or AI resources:
\begin{itemize}
\item TODO
\item TODO
\item $\ldots$	
\end{itemize}

By turning this document in, I attest that I have followed the 
\href{https://www.cs.tufts.edu/cs/136/2026f/index.html#collaboration}{[CS 136 Collaboration Policy]}. All solutions represent my own work. No solution text in this document was directly provided by other humans or artificial agents. 

\tableofcontents

\newpage
\officialdirections{
\subsection*{1a: Problem Statement}
Find the optimal one-hot assignment vectors $z^1$ for all $N=7$ examples, given the initial cluster locations $\mu^0$. Report $z^1$ and the value of the loss function $J(x, z^1, \mu^0)$.
}

\subsection{1a: Solution}
TODO FILL IN TABLE

\begin{tabular}{p{5cm} | p{5cm} | r}
$\mu^0$ & $z^1$ & $J(x_{1:N}, z^1, \mu^0)$
\\
\midrule
\begin{verbatim}
[[-3.  -2. ]
 [ 1.5  3. ]
 [ 2.   2. ]]
\end{verbatim}
&	
\begin{verbatim}
[[? ? ?]
 [? ? ?]
 [? ? ?]
 [? ? ?]
 [? ? ?]
 [? ? ?]
 [? ? ?]]
\end{verbatim}
&
 ??.?????
\end{tabular}


\newpage
\officialdirections{
\subsection*{1b: Problem Statement}
 Find the optimal cluster locations $\mu^1$ for all K=3 clusters, using the optimal assignments $z^1$ you found in 1a. Report $\mu^1$ and the value of the loss function $J(x, z^1, \mu^1)$.}

\subsection{1b: Solution}
TODO FILL IN TABLE

\begin{tabular}{p{5cm} | p{5cm} | r}
$\mu^1$ & $z^1$ & $J(x_{1:N}, z^1, \mu^1)$
\\
\midrule
\begin{verbatim}
[[ ?.???  ?.???]
 [ ?.???  ?.???]
 [ ?.???  ?.???]]
\end{verbatim}
&	
\begin{verbatim}
[[? ? ?]
 [? ? ?]
 [? ? ?]
 [? ? ?]
 [? ? ?]
 [? ? ?]
 [? ? ?]]
\end{verbatim}
&
 ??.?????
\end{tabular}


\newpage
\officialdirections{
\subsection*{1c: Problem Statement}
How does the loss function change from 1a to 1b? Does this align with known guarantees for the k-means algorithm?
}
\subsection{1c: Solution}
TODO YOUR SOLUTION HERE


\newpage
\officialdirections{
\subsection*{1d: Problem Statement}
Execute another iteration of k-means: use the $\mu^1$ from 1b to get new assignments $z^2$, and then use those assignments to update the locations $\mu^2$. 

Inspecting the values of $z^2$ and $\mu^2$, what interesting phenomenon do you see happening regarding the second cluster?
}
\subsection{1d: Solution}
TODO YOUR SOLUTION HERE


\newpage
\officialdirections{\subsection*{2a: Problem Statement}
The posterior $p(z_n|x_n)$ is a categorical distribution over the $K$-ary variable $z_n$. Show that, in the GMM setting defined here, its probabilities are given by
\begin{align*}
p(z_{nk}=1|x_n) = \frac
	{ \pi_k \text{exp}( - \frac{1}{2\epsilon} (x_n - \mu_k)^T(x_n - \mu_k) )}
	{ \sum_{j=1}^K \pi_j \text{exp}( - \frac{1}{2\epsilon} (x_n - \mu_j)^T (x_n - \mu_j) )}
\end{align*}
We call this probability $\gamma_{nk} := p(z_{nk}=1|x_n)$.
}

\subsection{2a: Solution}
TODO YOUR SOLUTION HERE

\newpage
\officialdirections{\subsection*{2b: Problem Statement}
Show that, if you treat $\gamma_{nk}$ as fixed values (which do not depend on $\mu_{1:K}$), the solution to 
\begin{equation*}
\argmax_{\mu_{1:K}}\,\,\, \sum_{n=1}^N  E_{p(z_{n}|x_n)}\left[\log p(z_n;\pi)p(x_n|z_n;\mu_{1:K},\epsilon) \right]
\end{equation*}
is given by
\begin{equation*}
\mu_k^* = \frac{1}{\sum_{n=1}^N \gamma_{nk}} \sum_{n=1}^N \gamma_{nk} x_n 
\end{equation*}
}

\subsection{2b: Solution}
TODO YOUR SOLUTION HERE

\newpage
\officialdirections{
\subsection*{3: Problem Statement}
Show that, if you fix $q(\theta_1)$ to some distribution, then the $q(\theta_2)$ which maximizes
\begin{equation*}
E_{q(\theta_1)q(\theta_2)}\left[\log \frac{p(\theta_1,\theta_2,x_{1:N})}{q(\theta_1)q(\theta_2)} \right]
\end{equation*}
out of the set of all possible distributions is given by
\begin{equation*}
\tilde{p}(\theta_2) = \frac{1}{C}\exp\Big\{E_{q(\theta_1)}\big[\log p(\theta_1,\theta_2,x_{1:N})\big]\Big\}
\end{equation*}
Where $C$ is the normalizing constant that makes $\tilde{p}(\theta_2)$ integrate to 1.
}

\subsection{3: Solution}

TODO YOUR SOLUTION HERE


\end{document}
